\documentclass[addpoints,10pt]{exam}

\usepackage{amsmath,amsthm,enumitem,wrapfig,amsfonts,mathtools}
\usepackage[mathscr]{euscript}
\usepackage[super]{nth}
\usepackage{dsfont}
\usepackage{xparse}
\usepackage{cancel}

\usepackage{geometry}
\usepackage[T1]{fontenc}
\renewcommand{\rmdefault}{ptm}
\usepackage{amsmath,amsfonts,amsthm,amssymb}
\usepackage{bm}
\usepackage{graphicx}
\usepackage{sectsty}
\usepackage{pgfplots}
\pgfplotsset{compat=1.18}
\usetikzlibrary{arrows.meta}
\usepackage{xcolor}
\definecolor{darkpastelgreen}{rgb}{0.01, 0.75, 0.24}
\definecolor{blue-violet}{rgb}{0.54, 0.17, 0.89}
\usepackage{polynom}

\newcounter{cprob}
\newenvironment{cprob}[1]{%
    \setcounter{cprob}{#1}%
    \noindent\textbf{Problem \thecprob.}%
}{%
    \par\bigskip%
}

\theoremstyle{plain}
\newcommand{\theoremname}{Theorem}
\newtheorem{thm}{\protect\theoremname}
\theoremstyle{definition}
\newtheorem{prob}[thm]{Problem}
\newtheorem*{problem*}{Open Problem}
\theoremstyle{plain}
\newtheorem{conjecture}[thm]{Conjecture}
\newtheorem{lem}[thm]{Lemma}
\newtheorem*{lem*}{Lemma}
\newtheorem{obs}[thm]{Observation}
\newtheorem{cor}[thm]{Corollary}
\theoremstyle{definition}
\newtheorem{definition}[thm]{Definition}

\let\oldprob\prob
\let\endoldprob\endprob
\renewenvironment{prob}
  {\begin{singlespace}\oldprob}
  {\endoldprob\end{singlespace}}

\newcommand{\belowtitle}{\leavevmode\newline}
\newcommand{\Observe}{\text{Observe.}}
\newcommand{\IF}{\mathbf{(\Rightarrow)}}
\newcommand{\FI}{\mathbf{(\Leftarrow)}}
\newcommand{\class}[2][]{\ensuremath{\left[\,#2\,\right]_{#1}}}

\newcommand{\horrule}[1]{\rule{\linewidth}{#1}}
\newcommand{\kkk}{\ensuremath{\Bbbk}}
\newcommand{\CC}{\ensuremath{\mathbb{C}}}
\newcommand{\FF}{\ensuremath{\mathbb{F}}}
\newcommand{\KK}{\ensuremath{\mathbb{K}}}
\newcommand{\NN}{\ensuremath{\mathbb{N}}}
\newcommand{\QQ}{\ensuremath{\mathbb{Q}}}
\newcommand{\RR}{\ensuremath{\mathbb{R}}}
\newcommand{\ZZ}{\ensuremath{\mathbb{Z}}}
\newcommand{\MM}{\ensuremath{\mathcal{M}}}
\newcommand{\TT}{\ensuremath{\mathcal{T}}}
\newcommand{\BB}{\ensuremath{\mathcal{B}}}
\newcommand{\VV}{\ensuremath{\mathcal{V}}}
\newcommand{\WW}{\ensuremath{\mathcal{W}}}
\newcommand{\UU}{\ensuremath{\mathcal{U}}}
\newcommand{\PP}{\ensuremath{\mathcal{P}}}
\newcommand{\LL}{\ensuremath{\mathcal{L}}}
\newcommand{\kk}{\ensuremath{\mathds{k}}}
\newcommand{\EE}{\ensuremath{\mathbb{E}}}
\newcommand{\Ss}{\ensuremath{\mathcal{S}}}
\newcommand{\GG}{\ensuremath{\mathcal{G}}}

\newcommand{\sm}{\char`\\}
\DeclarePairedDelimiter{\ip}{\langle}{\rangle}
\DeclarePairedDelimiter{\norm}{\lVert}{\rVert}
\DeclarePairedDelimiter{\sqb}{\lbrack}{\rbrack}

\newcommand{\floor}[1]{\left\lfloor #1 \right\rfloor}
\newcommand{\ceil}[1]{\left\lceil #1 \right\rceil}
\newcommand{\mbf}[1]{\ensuremath{\mathbf{#1}}}
\newcommand{\tbf}[1]{\textbf{ #1 }}
\newcommand{\Span}{\ensuremath{\mathrm{Span}}}
\newcommand{\Char}[1]{\mathrm{Char}\; #1}
\DeclareMathOperator{\lcm}{lcm}
\newcommand{\id}{\ensuremath{\mathrm{id}}}
\newcommand{\Gal}[2]{\mathrm{Gal}(#1/#2)}
\newcommand{\Aut}{\ensuremath{\mathrm{Aut}}}
\newcommand{\Fix}[2]{\mathrm{Fix}_{#1}(#2)}
\newcommand{\nil}{\ensuremath{\mathrm{nil}}}
\newcommand{\Hom}{\ensuremath{\mathrm{Hom}}}
\newcommand{\Obj}{\ensuremath{\mathrm{Obj}}}
\newcommand{\Rng}{\ensuremath{\mathrm{Rng}}}
\newcommand{\Ring}{\ensuremath{\mathrm{Ring}}}
\newcommand{\h}{\ensuremath{\mathrm{ht}}}

\makeatletter
\renewcommand*\env@matrix[1][*\c@MaxMatrixCols c]{%
  \hskip -\arraycolsep
  \let\@ifnextchar\new@ifnextchar
  \array{#1}}
\makeatother

\def\env@matrix{\hskip -\arraycolsep
  \let\@ifnextchar\new@ifnextchar
  \array{*\c@MaxMatrixCols c}}

\newcommand{\proj}[2]{\text{proj}_{#1}(#2)}

\newcommand{\StudentName}{Danny Banegas}
\newcommand{\AssignmentName}{Homework 5}
\newcommand{\CourseName}{MATH 720 - Algebra I}

\pagestyle{headandfoot}
\runningheadrule
\firstpageheadrule
\firstpageheader{\CourseName}{\StudentName}{\AssignmentName}
\runningheader{\CourseName}{\StudentName}{\AssignmentName}
\firstpagefooter{}{\thepage}{}
\runningfooter{}{\thepage}{}

\printanswers

\DeclareMathAlphabet{\mathcal}{OMS}{cmsy}{m}{n}

\usepackage{parskip}
\usepackage{setspace}

\begin{document}

%%%%%%%%%%%%%%%%%%%%%%%%%%%%%%%%%%%%%%%%%%%%%%%%% 1 %%%%%%%%%%%%%%%%%%%%%%%%%%%%%%%%%%%%%%%%
\setcounter{thm}{0}

\begin{prob}
A discrete valuation on a field $k$ is a map $v : k \to \ZZ \cup \{\infty\}$ such that $v(xy)=v(x)+v(y)$,
\[
v(x+y)\geq \min\{v(x),v(y)\},
\]
and
\[
v(x)=\infty \iff x=0.
\]
Prove that the set
\[
R:=\{a\in k\setminus\{0\}\mid v(a)\geq 0\}\cup\{0\}
\]
is a subring of $k$ and is in fact a Euclidean domain.

(Rings of this form are called discrete valuation rings or DVRs).
\end{prob}

\begin{proof}
For any $a,b\in R$,

$[1]:$
$v(1\cdot a)=v(1)+v(a)=v(a)\implies v(1)=0\implies 1\in R$.

$[+]:$
$v(a+b)\geq \min\{v(a),v(b)\}\geq 0$
since $a,b\in R\implies v(a),v(b)\geq 0\implies a+b\in R$.

$[\cdot]:$
$v(ab)=v(a)+v(b)\geq 0$
since $a,b\in R\implies ab\in R$.

$[-]:$
$v((-1)(-1))=v(1)=v(-1)+v(-1)=0$
so $v(-1)=0\implies -1\in R$ and then
$(-1)\cdot a=-a\in R$ by $[\cdot]$.

So $R\leq K$ is a subring.

Now consider any $a,b\in R$, $b\neq 0$. Since $K$ is a field,
$ba^{-1}\in K$, and
$v(ba^{-1})=v(b)+v(a^{-1})=v(b)-v(a)$.
Therefore, if $v(a)\geq v(b)$, then
$v(ab^{-1})=v(a)-v(b)\geq 0\implies ab^{-1}=\frac{a}{b}\in R$.

Also notice that $v(R)=\NN$. We will use this to prove $R$ is a
Euclidean domain.

For any $a,b\in R\setminus\{0\}$, either

(i) $v(a)\geq v(b)$, or

(ii) $v(a)<v(b)$.

(i) If $v(a)\geq v(b)$, let
$q=\frac{a}{b}\in R$, $r=0$. We have
$a=qb+r=\left(\frac{a}{b}\right)b+0=a$
with $r=0$.

(ii) If $v(a)<v(b)$, then let
$q=0\in R$ and $r=a\in R$. We have
$a=qb+r=(0)b+a$
and
$v(r)=v(a)<v(b)$.

So for all $a,b\in R\setminus\{0\}$, there exist $q,r\in R$ such that
$a=qb+r$ where $r=0$ or $v(r)<v(b)$ and $v(R)=\NN$. Thus $R$ is a
Euclidean domain.
\end{proof}

%%%%%%%%%%%%%%%%%%%%%%%%%%%%%%%%%%%%%%%%%%%%%%%%% 2 %%%%%%%%%%%%%%%%%%%%%%%%%%%%%%%%%%%%%%%%
\newpage
\begin{prob}
Let $I,J$ be ideals in a commutative ring $R$, and define the maps
\[
\varphi:R\to R/I\times R/J
\]
by
\[
r\mapsto ([r]_I,[r]_J)
\]
and
\[
\psi:R/I\times R/J\to R/(I+J)
\]
by
\[
([a]_I,[b]_J)\mapsto [a-b]_{I+J}.
\]
Thus,
\[
R\xrightarrow{\varphi}R/I\times R/J\xrightarrow{\psi}R/(I+J).
\]

\begin{enumerate}[label=(\alph*)]
\item Prove that $([a]_I,[b]_J)\in \operatorname{im}\varphi$ if and only if there exists some $x\in R$ such that
\[
[x]_I=[a]_I \qquad \text{and} \qquad [x]_J=[b]_J.
\]

\item Prove that $([a]_I,[b]_J)\in \ker\psi$ if and only if
\[
[a]_{I+J}=[b]_{I+J}.
\]
\end{enumerate}
\end{prob}

\begin{proof}
(a)
Suppose
$([a]_I,[b]_J)\in\varphi(R)$.
Then there exists $x\in R$ such that
$\varphi(x)=([a]_I,[b]_J)$, so
$([x]_I,[x]_J)=([a]_I,[b]_J)$.
So then
$[x]_I=[a]_I$
and
$[x]_J=[b]_J$.

Conversely, suppose there exists $x\in R$ such that
$[x]_I=[a]_I$
and
$[x]_J=[b]_J$.
Then
$\varphi(x)=([x]_I,[x]_J)=([a]_I,[b]_J)$,
so
$([a]_I,[b]_J)\in\varphi(R)$.

(b)
Suppose
$([a]_I,[b]_J)\in\ker\psi$.
Then
$\psi(([a]_I,[b]_J))=[0]_{I+J}$,
so
$[a-b]_{I+J}=[0]_{I+J}$.
Hence
$a-b\in I+J$,
which is equivalent to
$[a]_{I+J}=[b]_{I+J}$.

Conversely, suppose
$[a]_{I+J}=[b]_{I+J}$.
Then
$[a-b]_{I+J}=[0]_{I+J}$,
so
$\psi(([a]_I,[b]_J))=[0]_{I+J}$.
Therefore
$([a]_I,[b]_J)\in\ker\psi$.

\end{proof}

%%%%%%%%%%%%%%%%%%%%%%%%%%%%%%%%%%%%%%%%%%%%%%%%% 3 %%%%%%%%%%%%%%%%%%%%%%%%%%%%%%%%%%%%%%%%
\newpage
\begin{prob}
Using the same notation as Problem 2:

\begin{enumerate}[label=(\alph*)]
\item Prove that
\[
\operatorname{im}\varphi=\ker\psi.
\]
In this case, we say the sequence is exact.

Note that when $R=\ZZ$, this becomes the statement:
\[
x\equiv a \pmod n
\qquad \text{and} \qquad
x\equiv b \pmod m
\]
has a solution if and only if
\[
a\equiv b \pmod{\operatorname{lcm}(m,n)}.
\]

\item Find integers $a,b,m,n$ such that
\[
\gcd(m,n)\neq 1
\]
but the system
\[
x\equiv a \pmod n,
\qquad
x\equiv b \pmod m
\]
still has a solution.
\end{enumerate}
\end{prob}

\begin{proof}
(a)
Suppose
$([a]_I,[b]_J)\in\varphi(R)$.
By the previous problem, there exists $x\in R$ such that
$[x]_I=[a]_I$
and
$[x]_J=[b]_J$.
Hence
$a-x\in I$
and
$b-x\in J$.
Therefore
$a-b=(a-x)+(x-b)\in I+J$,
so
$[a]_{I+J}=[b]_{I+J}$.
By the previous problem again,
$([a]_I,[b]_J)\in\ker\psi$.

Conversely, suppose
$([a]_I,[b]_J)\in\ker\psi$.
Then
$[a]_{I+J}=[b]_{I+J}$,
so
$a-b=i+j$
for some $i\in I$ and $j\in J$.

Let
$x=a-i=b+j$.
Then
$x-a=-i\in I$
and
$x-b=j\in J$,
so
$[x]_I=[a]_I$
and
$[x]_J=[b]_J$.
Hence
$([a]_I,[b]_J)\in\varphi(R)$.

Therefore
$\varphi(R)=\ker\psi$.

(b)
Let $a=b=0$ and $m=2,n=4$. $(a)\implies \begin{cases} x\equiv a\pmod m\\ x\equiv a\pmod m\end{cases}\iff a\equiv b\pmod \lcm(m,n)=2.$ This system has the solution $x=8$ and $\gcd(2,4)=2\neq 1$.

\end{proof}

%%%%%%%%%%%%%%%%%%%%%%%%%%%%%%%%%%%%%%%%%%%%%%%%% 4 %%%%%%%%%%%%%%%%%%%%%%%%%%%%%%%%%%%%%%%%
\newpage
\begin{prob}
Review the notes on Euclidean domains, PIDs, and UFDs (Sections 1.13--1.15).

\begin{enumerate}[label=(\alph*)]
\item If $R$ is a Euclidean domain and $I\subseteq R$, then we proved $I$ is principal, say
\[
I=(a)
\]
for some $a\in R$. Describe $a$.

\item Outline the proof that every PID is a UFD. You do not need to reproduce the full proof, but list every major step.
\end{enumerate}
\end{prob}

\begin{proof}
(a)
Let $R$ be a Euclidean domain with Euclidean valuation
$\delta:R\setminus\{0\}\to\NN$, and let $I\subseteq R$ be a nonzero ideal.

Since $\delta(I\setminus\{0\})\subseteq\NN$, by the \textbf{Well-Ordering Principle} there exists $a\in I\setminus\{0\}$ such that $\delta(a)$ is
minimal.

We prove that $I=\langle a\rangle$. Obviously, $\langle a\rangle\subseteq I$ since $a\in I.$

Now let $x\in I$. Since $R$ is Euclidean, there exist $q,r\in R$ such
that
$x=qa+r$,
where either $r=0$ or $\delta(r)<\delta(a)$.

Since $x,a\in I$ and $I$ is an ideal,
$r=x-qa\in I$.

By minimality of $\delta(a)$, we can't have that $\delta(r)<\delta(a)$. So $r=0$.

Therefore $x=qa\in\langle a\rangle$, so
$I\subseteq\langle a\rangle$.

So each $I=\langle a\rangle$ where $a$ minimizes the Euclidean evaluation $\delta$ among non-zero elements of $I$.

(b)
Let $R$ be a PID.

We first prove existence of irreducible factorizations. Suppose there
exists a nonzero non-unit which does not factor into irreducibles.
Choose such an element $x_1$.

Since $x_1$ is not irreducible, we may write
$x_1=x_2y_2$,
where neither $x_2$ nor $y_2$ is a unit. In particular,
$x_1\in\langle x_2\rangle$, so
$\langle x_1\rangle\subseteq\langle x_2\rangle$.

If equality held, then
$x_2\in\langle x_1\rangle$, so
$x_2=rx_1$
for some $r\in R$. Hence
\[
x_1=x_2y_2=(rx_1)y_2=x_1(ry_2).
\]
Since $R$ is an integral domain and $x_1\neq 0$, we get
$1=ry_2$, so $y_2$ is a unit, a contradiction.

Thus
$\langle x_1\rangle\subsetneq\langle x_2\rangle$.

Continuing inductively gives a strictly increasing chain
\[
\langle x_1\rangle
\subsetneq
\langle x_2\rangle
\subsetneq
\langle x_3\rangle
\subsetneq
\cdots.
\]

But every PID is Noetherian, so such a chain cannot exist and every
nonzero non-unit factors into irreducibles.

Now let $p$ be irreducible and suppose $p\mid ab$. Consider the ideal
$\langle p,a\rangle$. Since $R$ is a PID,
$\langle p,a\rangle=\langle d\rangle$
for some $d\in R$.

Since $d\mid p$ and $p$ is irreducible, either $d$ is a unit or $d$ is
associate to $p$.

If $d$ is associate to $p$, then
$\langle p,a\rangle=\langle p\rangle$,
so $a\in\langle p\rangle$, i.e. $p\mid a$.

If $d$ is a unit, then
$\langle p,a\rangle=R$, so there exist $x,y\in R$ such that
$1=xp+ya$.
Multiplying by $b$ gives
$b=xpb+yab$.

Since $p\mid ab$ and clearly $p\mid xpb$, it follows that $p\mid b$.

Therefore every irreducible is prime.

Finally, suppose
$p_1\cdots p_n=q_1\cdots q_m$
with all $p_i,q_j$ irreducible. Since $p_1$ is prime,
$p_1\mid q_1\cdots q_m$, so $p_1\mid q_j$ for some $j$. Since $q_j$ is
irreducible, $p_1$ is associate to $q_j$. After reordering and cancelling
associates, continue inductively and you can find some associate $q_{j}$ for each $p_{i}$.

So we see factorization is unique up to associates and reordering, and $R$
is a UFD.
\end{proof}

%%%%%%%%%%%%%%%%%%%%%%%%%%%%%%%%%%%%%%%%%%%%%%%%% 5 %%%%%%%%%%%%%%%%%%%%%%%%%%%%%%%%%%%%%%%%
\newpage
\begin{prob}
Let $R$ be an integral domain, and let $p$ be a prime element of $R$.

\begin{enumerate}[label=(\alph*)]
\item Prove that $p$ is not a unit.

\item Prove that if $q$ is an associate of $p$, then $q$ is also prime.
\end{enumerate}
\end{prob}

\begin{proof}
(a)
Suppose $p$ is a unit. Then
$\langle p\rangle=\langle 1\rangle=R$,
which is not a proper ideal. Since a prime ideal must be proper, this is
a contradiction. Therefore $p$ is not a unit.

(b)
Suppose $q$ is associate to $p$. Then
$q=up$
for some unit $u\in R^\times$.

Now suppose $q\mid ab$. Then
$ab=qk=upk$
for some $k\in R$, so $p\mid ab$.

Since $p$ is prime, either $p\mid a$ or $p\mid b$.

If $p\mid a$, then
$a=pr$
for some $r\in R$. Since $p=u^{-1}q$, we have
\[
a=pr=(u^{-1}q)r=q(u^{-1}r),
\]
so $q\mid a$.

Similarly, if $p\mid b$, then $q\mid b$.

Therefore $q$ is prime.
\end{proof}

%%%%%%%%%%%%%%%%%%%%%%%%%%%%%%%%%%%%%%%%%%%%%%%%% 6 %%%%%%%%%%%%%%%%%%%%%%%%%%%%%%%%%%%%%%%%
\newpage
\begin{prob}
Let $R$ be a commutative ring and let $S\subseteq R$ be a multiplicative subset of $R$.

\begin{enumerate}[label=(\alph*)]
\item Prove that if $I$ is an ideal of $R$ such that
\[
I\cap S=\emptyset,
\]
then
\[
I^e:=S^{-1}I=\left\{\frac{a}{b}\mid a\in I,\ b\in S\right\}
\]
is a proper ideal of $S^{-1}R$.

\item Let
\[
\varphi:R\to S^{-1}R
\]
be the natural ring homomorphism defined by
\[
r\mapsto \frac r1.
\]
Prove that if $J$ is a proper ideal of $S^{-1}R$, then
\[
J^c:=\varphi^{-1}(J)
\]
is an ideal of $R$ such that
\[
J^c\cap S=\emptyset.
\]
\end{enumerate}
\end{prob}

\begin{proof}
(a)
$\forall \frac{i_1}{s_1},\frac{i_2}{s_2}\in S^{-1}I,\text{ we have that }i_1,i_2\in I$ and $s_1,s_2\in S$. So then 
$s_2i_1-s_1i_2\in I$ and
\[
\frac{i_1}{s_1}+\frac{i_2}{s_2}
=
\frac{s_2i_1+s_1i_2}{s_1s_2}
\in S^{-1}I\leq_{+}S^{-1}R.
\]

$\forall\frac{r}{s}\in S^{-1}R$ we have
\[
\frac{r}{s}\cdot\frac{i_1}{s_1}
=
\frac{ri_1}{ss_1}\text{ and since $I$ is an ideal, $ri_1\in I$, so
$\frac{ri_1}{ss_1}\in S^{-1}I$.}
\]


Therefore $S^{-1}I$ is an ideal of $S^{-1}R$.

(b)
Since $\varphi:R\to S^{-1}R$ is a ring homomorphism and $J$ is a proper ideal
$J^c=\varphi^{-1}(J)$ is an ideal of $R$. Now suppose $s\in J^c\cap S$. Then
$\varphi(s)=\frac{s}{1}\in J$.

But
\[
\frac{s}{1}\cdot\frac{1}{s}=\frac{1}{1},
\]
so $\frac{s}{1}$ is a unit in $S^{-1}R$. Since $\frac{s}{1}\in J$ and $\frac{1}{s}\in S^{-1}R$
we get $\frac{1}{1}=1_{S^{-1}R}\in J$ and so $J=S^{-1}R$, contradicting that $J$ is proper.

Thus $J^c\cap S=\emptyset$.

\end{proof}

%%%%%%%%%%%%%%%%%%%%%%%%%%%%%%%%%%%%%%%%%%%%%%%%% 7 %%%%%%%%%%%%%%%%%%%%%%%%%%%%%%%%%%%%%%%%
\newpage
\begin{prob}
Using the same notation as Problem 6, let $J$ be a proper ideal of $S^{-1}R$.

\begin{enumerate}[label=(\alph*)]
\item Prove that
\[
(J^c)^e=J
\]
while
\[
(I^e)^c=\{a\in R\mid (\exists s\in S)\text{ such that } sa\in I\}.
\]

\item Find an example showing that $(I^e)^c$ need not equal $I$, even if $I\cap S=\emptyset$.

(Hint: Let
\[
S=\{1,x,x^2,\dots\}
\]
in $R=\CC[x,y]$, and consider $I=(xy)$.)
\end{enumerate}
\end{prob}

\begin{proof}
For Problem 7:

(a)
First we prove $(J^c)^e=J$. Since $J^c=\varphi^{-1}(J)$, every
$a\in J^c$ satisfies $a/1\in J$. So every element of $(J^c)^e$ is of
the form $(r/s)(a/1)$ with $a/1\in J$ and $(J^c)^e\subseteq J$.

Conversely, let $a/s\in J$. Since $J$ is an ideal and $s/1\in S^{-1}R$,
we have $(s/1)(a/s)=a/1\in J$. So $a\in J^c$, and
$a/s=(1/s)(a/1)\in (J^c)^e$. Therefore $J\subseteq (J^c)^e$, and
$(J^c)^e=J$.

Now we prove
$(I^e)^c=\{a\in R\mid (\exists s\in S)\text{ such that }sa\in I\}$.

Let $a\in (I^e)^c$. Then $a/1\in I^e=S^{-1}I$. Hence
$a/1=i/t$ for some $i\in I$ and $t\in S$. By equality in the
localization, there exists $u\in S$ such that $u(ta-i)=0$. Thus
$uta=ui\in I$. Since $ut\in S$, this shows there exists $s\in S$ such
that $sa\in I$.

Conversely, suppose there exists $s\in S$ such that $sa\in I$. Then
$a/1=sa/s\in S^{-1}I=I^e$ and $a\in (I^e)^c$. Therefore
$(I^e)^c=\{a\in R\mid (\exists s\in S)\text{ such that }sa\in I\}$.

(b)
Let $R=\CC[x,y]$, let $S=\{1,x,x^2,\dots\}$, and let $I=\langle xy\rangle$.
Then $I\cap S=\emptyset$, since no power of $x$ is divisible by $xy$.

However, $y\in (I^e)^c$, because $xy\in I$ and $x\in S$. But
$y\notin I=\langle xy\rangle $. Therefore, $(I^e)^c\neq I$.
\end{proof}

%%%%%%%%%%%%%%%%%%%%%%%%%%%%%%%%%%%%%%%%%%%%%%%%% 8 %%%%%%%%%%%%%%%%%%%%%%%%%%%%%%%%%%%%%%%%
\newpage
\begin{prob}
Let $R$ be a UFD, and let $S$ be a multiplicative subset of $R$.

\begin{enumerate}[label=(\alph*)]
\item Prove that if $q$ is irreducible in $R$, then $q/1$ is either irreducible or a unit in $S^{-1}R$.

\item Prove that if $a/s$ is irreducible in $S^{-1}R$, then $a/s$ is associate to $q/1$ for some irreducible element $q$ of $R$.

\item Prove that $S^{-1}R$ is also a UFD.
\end{enumerate}
\end{prob}

\begin{proof}
(a)
Let $q$ be irreducible in $R$. Since $R$ is a UFD, $q$ is prime.

If $q/1$ is a unit in $S^{-1}R$, then we are done. So suppose $q/1$ is
not a unit. We prove $q/1$ is irreducible.

Suppose $q/1=(a/s)(b/t)$ in $S^{-1}R$. Then, after clearing
denominators, there exists $u\in S$ such that $uqst=uab$. Since $R$ is
a domain, $qst=ab$. Since $q$ is prime, $q\mid a$ or $q\mid b$.

If $q\mid a$, we have $a=qa'$ for some $a'\in R$. Then
$q/1=(qa'/s)(b/t)$, and then cancelling $q/1$ gives
$1=(a'/s)(b/t)$. So $b/t$ is a unit. Similarly, if $q\mid b$, then
$a/s$ is a unit. Therefore $q/1$ is irreducible in $S^{-1}R$.

(b)
Suppose $a/s$ is irreducible in $S^{-1}R$. Factor $a$ in the UFD $R$ as
$a=u q_1\cdots q_n$, where $u\in R^\times$ and each $q_i$ is irreducible
in $R$.

Then in $S^{-1}R$,
$a/s=(u/1)(q_1/1)\cdots(q_n/1)(1/s)$.
Here $u/1$ and $1/s$ are units. Since $a/s$ is irreducible, exactly one
of the factors $q_i/1$ is nonunit, and all the others are units.
Therefore $a/s$ is associate to $q_i/1$ for some irreducible $q_i\in R$.

(c)
Every nonzero nonunit $a/s\in S^{-1}R$ factors into irreducibles by
factoring $a$ in $R$ and discarding the factors which become units in
$S^{-1}R$.

By part (b), every irreducible in
$S^{-1}R$ is associate to $q/1$ for some irreducible $q\in R$.

Also, if $q/1$ is irreducible in $S^{-1}R$, then it is prime. Suppose $(q/1)\mid (a/s)(b/t)$. Clearing denominators gives
$wuab=wqcst$ for some $w,u\in S$. Since $q$ is prime in $R$, $q$ divides
one of $w,u,a,b,s,t$. But if $q$ divided one of $w,u,s,t\in S$, then
$q/1$ would be a unit in $S^{-1}R$, contradiction. So $q\mid a$ or
$q\mid b$, and then $q/1$ divides $a/s$ or $b/t$.

Thus irreducibles in $S^{-1}R$ are prime. The usual prime-divisibility
argument gives uniqueness of factorizations up to associates and order.
Therefore $S^{-1}R$ is a UFD.

\end{proof}

\end{document}
